% =====================================================================
% SECTION V — EXPERIMENTAL SETUP
% File: section5_setup.tex
% =====================================================================

% ------------------------------------------------------------------
\section{Experimental Setup}
\label{sec:setup}
% ------------------------------------------------------------------

% ------------------------------------------------------------------
\subsection{Datasets}
% ------------------------------------------------------------------

\textbf{SupplyGraph}~\cite{wasi2024supplygraph} is a supply-chain
demand dataset containing daily operational demand for 40 products
across 5 EEA product groups and 222 trading days. The graph encodes
4 typed edge relations (product group, sub-group, plant, storage
location) over 40 nodes. The dataset exhibits pronounced
zero-inflation and intra-group demand sparsity. A strict
chronological 70/10/20 split is applied; because of the 222-day
window, the held-out test partition is populated predominantly by
the EEA group. We report this distributional artefact transparently
and confirm that validation loss is computed across the full 40-node
graph.

\textbf{USGS Mineral Resources}~\cite{usgs2023minerals} is derived
from the USGS National Minerals Information Center annual
production and import statistics for 127 mineral commodities over
five years (2019--2023). The graph encodes 15 typed edge relations
(supply-chain input, technology cluster, critical mineral
designation, geopolitical supply risk, and eleven additional
process and application linkages). The annual resolution means that
each seed run trains on only 3--4 years of data, posing a
low-data-regime challenge distinct from SupplyGraph.

\textbf{Lead-time availability.} Neither dataset provides operational transport or procurement lead-time logs. To enable functional evaluation of the downstream risk-scoring interface (Layer 5), we assign proxy lead times: for SupplyGraph, group-level constants sampled once per group from a bounded range [3, 30] days; for USGS, a fixed constant of 1.0 reflecting the absence of shipment granularity in annual aggregates. These proxies are used solely to verify the computational integration of the risk layer and do not represent real-world logistics conditions. Consequently, risk-related outputs should be interpreted as illustrative rather than operationally validated. The forecasting model itself is independent of these proxies; lead times affect only the post-hoc risk layer.

% ------------------------------------------------------------------
\subsection{Training Protocol}
% ------------------------------------------------------------------

All models are trained for 70 epochs with the Adam
optimiser~\cite{kingma2015adam}, initial learning
rate $1{\times}10^{-3}$, and a step-decay schedule (halved at
epochs 27, 47, 63). Batch normalisation is replaced by LayerNorm
throughout. Dropout is 0.3 (USGS) and 0.2 (SupplyGraph).
To quantify reproducibility, every experiment is repeated across
five random seeds (42, 123, 456, 789, 1337); all reported metrics
are mean $\pm$ standard deviation across seeds. Training runs on
a single NVIDIA T4 GPU via Google Cloud Platform (Vertex AI
Workbench); model code uses PyTorch~\cite{paszke2019pytorch} and
PyTorch Geometric~\cite{fey2019pyg}.

\smallskip
\noindent\textbf{Disclosure.} Computational cost data (GPU-hours,
CO\textsubscript{2} impact) were not systematically recorded and
are not reported. All five seeds share identical architecture and
hyperparameters; no seed-specific tuning was performed.

% ------------------------------------------------------------------
\subsection{Baselines}
% ------------------------------------------------------------------

We compare against six models:

\begin{itemize}[leftmargin=*, nosep]
  \item \textbf{ARIMA} --- classical univariate autoregressive model;
    serves as a statistical lower bound~\cite{hyndman2018fpp}.
  \item \textbf{XGBoost} — gradient-boosted trees with lag features
    \cite{chen2016xgboost}; the strongest classical baseline.
  \item \textbf{TFT-only} — our temporal branch alone, without the
    GAT branch or adaptive fusion.
  \item \textbf{GAT-only} — our structural branch alone.
  \item \textbf{Fixed Fusion} ($\alpha{=}0.5$) — equal, non-learned
    weighting of both branches; ablates the adaptive fusion
    mechanism specifically.
  \item \textbf{Adaptive Fusion (Ours)} — full model with learned
    $\alpha$ conditioned on material type and horizon.
\end{itemize}

Classical baselines use the same chronological split and are fitted
independently per seed for fair comparison. ARIMA and XGBoost do
not use graph structure; this disadvantage is inherent to the
comparison and is not corrected for.

% ------------------------------------------------------------------
\subsection{Evaluation Metrics}
% ------------------------------------------------------------------

We report WAPE (primary), RMSE, SMAPE, R$^2$, P10--P90 coverage,
mean interval width, and P50 WAPE. WAPE is the primary metric
because it is scale-invariant and robust to zero-inflation
~\cite{makridakis2022m5}. Statistical significance is assessed with
the one-tailed Wilcoxon signed-rank test~\cite{wilcoxon1945} across
the five seed outcomes; $p{=}0.0625$ is the minimum achievable
p-value for $n{=}5$~\cite{gneiting2007scoring}.
